\Chapter{49}

\section*{The Blessing of Jacob}

We take up a distinctly different genre of ancient Yahwistic poetry with the
tribal blessings, the Blessing of Jacob and the Blessing of Moses. To the same
genre belong the tribal references in Jud. 5:14--18 and Ps. 68:28.\textsuperscript{a}
The Blessings also have certain characteristics in common with the Oracles of
Balaam, but there are also strong contrasts.\textsuperscript{b}

It seems probable in the light of \emph{gattungs- und überlieferungs-geschichtlich}
studies that individual blessings rest on ancient tradition. They circulated in
oral form as folk literature, finally being gathered into collections. Some
blessings are composite, such as those of Judah (Gen. 49:8--12) and Levi
(Deut. 33:8--11), made up of originally separate elements. These illustrate the
way in which blessings were brought together. On the other hand, the tendency
to chop up the blessings into a host of fragments is the result of overly
atomistic analysis. The Blessing of Jacob consists of a nucleus of blessings
consciously set in the mouth of the Patriarch (e.g., Reuben, Judah), to which
have been appended blessings drawn from other sources (e.g., the Joseph
blessing, which occurs in slightly altered form in the Blessing of Moses). We
may assume that groups of blessings, ascribed to Jacob and Moses, and perhaps
others, circulated orally in the period of the Judges.

After the earlier classical pieces, the Songs of Deborah and Miriam, the
blessings in these two collections are disappointing as literature. They are
practically folk poetry, and lack the varied meter and complex rhetorical style
of the victory odes. They delight in plays on words, frequently use rhyme, and
with certain notable exceptions are uninspired in content.

With our increasing knowledge of the amphictyonic organization of the tribes
in the period of the Judges,\textsuperscript{c} the place of the collection of
tribal blessings in the religious and political life of the people has become
clear. One may be assured that such collections were framed in this period from
various cycles of oral tradition. It is most unlikely that the blessings were
artificially set in a framework after the ``twelve''-tribe system had ceased to
be a living organization.

No serious question has been raised by scholars as to the pre-monarchic date of
the majority of the individual blessings in Gen. 49. The only exception worthy
of comment is the blessing of Judah, Gen. 49:8--11. As will be pointed out in
the study which follows, this blessing is composite: vss. 8, 9, and 11--12 are
individual units of a most archaic type. The short enigmatic blessing, vs. 10,
is usually referred to the Judahite monarchy, the earliest possibility being
the reign of David. An allusion to David is possible; more probably we have
here pre-Davidic material which came naturally to be associated with the great
king, was interpreted and may have been slightly modified in the light of his
career and achievements. This is speculative, however, and it is quite
impossible to date the collection on the basis of a doubtful understanding of
a corrupt text. Judahite hegemony in the south during the period of the Judges
may be all that is implied. It may be remarked that certain blessings (e.g.,
Reuben, Simeon and Levi) may go back to patriarchal traditions. It is best to
date the collection in substantially its present form in the late period of the
Judges. Orthographic and linguistic criteria support such a date.


\Verse{1b}
{
	האסף ואגד לכם [ ]
}
{
	Gather yourselves that I may inform you,
}
{
}


\Verse{2}
{
	הקבצו [ ] בני יעקב

	ושמעו אל ישראל אבכם
}
{
	Assemble together, O sons of Jacob,

	Hearken unto Israel your father.
}
{
}


\Verse{3}
{
	ראובן בכר

	אתה כח

	וראשית אן

	יתר עז
}
{
	Reuben, my first-born

	Thou art my strength

	The prime of my manhood

	The best of my might.
}
{
}


\Verse{4}
{
	כ עלת <על> משכבי

	אז חללת יצועי <אבכם>
}
{
	For thou didst go up to my bed

	Yea, thou didst profane the couch
	of thy father.
}
{
}


\Verse{5}
{
	שמעון ולוי אחים

	כלי חמס מכרתם [ ]
}
{
	Simeon and Levi are brethren

	Weapons of violence are their merchandise.
}
{
}


\Verse{6}
{
	בסדם אל תבא נפש

	בקהלם אל תחד כבדי

	[ ] באפם הרג איש

	וברצנם עקר שור
}
{
	Into their council enter not, O my soul

	In their assembly join not, O my heart.

	In their anger, they slew a man

	Willfully they houghed a bull.
}
{
}


\Verse{7}
{
	ארר אפם כי עז

	ועברתם כי קשת

	אחלקם ביעקב

	אפיצם בישראל
}
{
	Cursed be their anger, for it is fierce

	And their wrath, for it is harsh.

	I will divide them in Jacob,

	I will scatter them in Israel.
}
{
}


\Verse{8}
{
	יהודה הוד אתה

	ידך בערף איביך

	יהודה ידך אחיך

	ישתחוו לך בני אביך
}
{
	Judah, my splendor art thou,

	Thy hand is on the neck of thine enemies.

	Judah, thy brothers shall praise thee

	The sons of thy father shall bow
	down to thee.
}
{
}


\Verse{9}
{
	גור ארי יהד

	מטרף בני עלת

	כרע רבץ כארי

	כלביא מי יקימנו
}
{
	A lion's whelp is Judah

	From the prey, my son, thou hast gone up.

	He crouches, he couches like a lion,

	Like a lioness, who dares rouse him?
}
{
}


\Verse{10}
{
	לא יסר שפט מיהד

	ומחקק מבין רגלו

	\_\_\_\_\_\_

	\_\_\_\_\_\_
}
{
	There shall not fail a judge from Judah

	Nor a commander from
	among his standards

	\ldots\ldots\ldots

	\ldots\ldots\ldots
}
{
}


\Verse{11}
{
	אסר לגפן ערו

	ולשרק בני אתנו

	כבס בין לבשו

	ובדם ענבים סותו
}
{
	His ass is harnessed with its bridle (?)

	The colt of his she-ass with its trappings (?)

	He dyes with wine his garment

	With the blood of grapes his mantle.
}
{
}


\Verse{12}
{
	חכללי עינם מיין

	ולבן שנים מחלב
}
{
	Darker are his eyes than wine

	Whiter his teeth than milk.
}
{
}


\Verse{13}
{
	זבלון

	לחף ימים ישב

	על מפרצ[ו] ישכן

	\_\_\_\_\_\_

	\_\_\_\_\_\_
}
{
	Zebulun:

	At the shore of the seas he settles,

	At its inlets, he encamps.

	\ldots\ldots\ldots

	\ldots\ldots\ldots
}
{
}


\Verse{14}
{
	יששכר חמר \_\_\_
	
	רבץ בין המשפתים
}
{
	Issachar is an ass \ldots

	He couches among the rubbish piles
}
{
}


\Verse{15}
{
	וירא מנחת כי טוב

	ויבט ארץ כי נעם

	ויט שכמו לסבל

	ויהי למס עבד
}
{
	He sees for himself a resting place
	which is good

	He beholds for himself a land
	which is pleasant

	He bends his shoulder to bear burdens

	And he has become a corvée slave.
}
{
}


\Verse{16}
{
	דן ידין <דין> עם

	כאחד שבטי ישראל
}
{
	Dan pleads the cause of his people,

	The tribes of Israel together.
}
{
}


\Verse{17}
{
	[ ] דן נחש על דרך

	שפיפן על ארח

	הנשך עקבי סס

	ויפל רכבו אחר [ ]
}
{
	Dan is a serpent on the road,

	A venomous snake on the path

	Who snaps at the heels of a horse,

	So that its rider falls backward.
}
{
}

